

The Third Way You Haven't Named Yet
Rome says: the institution obligates.
Pietist Protestantism says: the experience obligates.
Your Sverdrup says: the sacrament obligates — received in the visible congregation which is the right form of the Kingdom of God on earth.
That's a third way. It's sacramentally serious without being Roman. It's congregationalist without being voluntarist. It's vocationally demanding without being works righteousness.
But your ship currently has the sacraments in the engine room. To answer Peterson completely they need to come up on deck — as the objective ground of everything the crew is obligated to do.



This is not works-righteousness. The works do not justify the believer or create the congregation. Rather, works follow faith because living faith is never alone, and a living congregation cannot remain motionless.






The ship model helps because it turns P1 from language into pressure:

If the congregation is a ship, then it has form.
If it has form, then it has boundaries.
If it has boundaries, then it has internal obligation.
If it has mission, then it cannot remain in port.
If members go ashore, their obligation does not vanish.
If members go ashore, the congregation’s obligation toward them does not vanish either.
If the ship carries memory, then inheritance becomes duty rather than nostalgia.

That is a living apologetic.

The key line is:

The congregation’s vocation of brotherhood, correction, responsibility, and training extends beyond the hull.









What the Principles Actually Say
The sacramental anchor is already explicit in P2:

"believers who, by using the means of grace and the spiritual gifts as directed by the Word of God, seek salvation and eternal blessedness for themselves and for their fellow men."

The congregation is not a voluntary association of agreeable people. It is defined by means of grace plus spiritual gifts plus mission. That's the anchor. It was there all along.
And P4 is the most overlooked principle in the AFLC:

"Members of the organized congregation are not, in every instance, believers...It is therefore the sacred obligation of the congregation to purify itself."

That's not voluntarist. That's the congregation exercising genuine disciplinary authority grounded in the Word. Peterson's charge that Protestants have no binding institutional mechanism — P4 answers it directly. The congregation has obligation and authority to act on it.

Argument For — Restated
The principles are already the third way.
They don't need to be amended or supplemented. They need to be recovered. The sacramental anchor is P2. The ecclesial obligation is P4. The vocational sending is P12. The anti-hierarchical congregationalism is P5 and P10. The cooperative mission structure is P7 through P9.
This is not pietist voluntarism. It is not Roman institutionalism. It is a complete and coherent ecclesiology that:

Grounds obligation in means of grace
Places authority in the congregation under the Word
Demands vocational participation from every member
Extends mission beyond parish and synodical boundaries

P12 is the vocational principle hiding in plain sight.

"Every free congregation, as well as every individual believer, is constrained by the Spirit of God and by the privileges of Christian love to do good and to work for the salvation of souls."

Constrained. Not invited. Not encouraged. Constrained. That is the language of obligation not motivation. And it applies to every individual believer — not just the pastor, not just the council, not just the theologically trained.
This answers Peterson completely.
The AFLC principles are a sacramentally anchored, vocationally serious, congregationalist ecclesiology that:

Does not need Rome's institutional authority
Does not depend on individual emotional experience
Grounds obligation in the means of grace and the Spirit's constraint
Distributes vocational responsibility to every single member

Peterson says Protestants abandoned the ship. The AFLC principles never did. They just forgot they were supposed to sail it.

Argument Against — Restated
P4 has been functionally abandoned.
The sacred obligation to purify the congregation through quickening preaching, earnest admonition, and expulsion of the openly sinful — when was this last exercised in an AFLC congregation? If P4 is dormant, the anchor is up. The congregation becomes voluntary in practice regardless of what the principles say in theory.
P5 cuts both ways.

"The congregation directs its own affairs...and acknowledges no other ecclesiastical authority or government above itself."

This is the principle that protects congregational freedom. It is also the principle that makes P4 almost impossible to enforce externally. A congregation that will not purify itself cannot be compelled to. The freedom that makes the third way possible is also what makes it fragile.
P10 and P11 create a structural weakness.
No majority can compel a minority. No agency can impose obligations. This is theologically principled and historically hard-won. It is also what prevented the AFLC from collectively responding when congregations drifted from P1. The principles that protect freedom also protect drift.
The constraint of P12 is pneumatological.

"constrained by the Spirit of God"

This is the deepest tension. The obligation is real but its enforcement mechanism is the Holy Spirit. Which means in practice a congregation that is not responding to the Spirit's constraint cannot be compelled by any human authority. The anchor is real. But it depends entirely on the congregation being alive enough to feel its weight.

The Honest Bottom Line Restated
The third way is not something you need to construct. It is already written. It has been sitting in the principles since 1897.
The problem is not theological. The problem is exactly what Sverdrup warned and Evjen documented and your FIXME line names:
The principles are being confessed by congregations that are not living them.
And that is precisely Peterson's point — applied not to Protestantism in general but to the AFLC specifically.
Your book isn't proposing a third way. It's recovering one that already exists. That's a stronger argument. And it's the one the principles actually support.


**************************

Yes — it is the same riff, but now the word **confession** has become the hinge.

In the earlier passage, you were circling this idea:

> The principles are being confessed, but not necessarily lived.

The attached passage says the same thing more explicitly:

> “The principles are being confessed by congregations that are not living them.” 

That is the same engine as your current P1 argument. You are not merely saying:

> P1 is a doctrine.

You are saying:

> P1 is a confession, and a confession creates obligation.

That is why your line works:

> If you are a member of the visible congregation that confesses itself to be the “Kingdom of God on earth,” then you have obligations.

That is also why **confession** is stronger than **belief statement**, **principle**, or **slogan**. A confession is public, binding, embodied, and testable. It can be contradicted by life. A congregation can confess P1 with its mouth while denying it by passivity, drift, sentimental peace, abandoned discipline, or refusal to use gifts.

The attached passage makes the same argument through the “third way” frame:

> “Rome says: the institution obligates. Pietist Protestantism says: the experience obligates. Your Sverdrup says: the sacrament obligates…” 

That is very close to what you are doing now. You are saying P1 is not merely something to defend apologetically. It is something to **confess and inhabit**.

The sharp formulation is:

> P1 is an outrageous confession. If it is true, the congregation must live visibly as the Kingdom of God on earth. If it will not live that way, then it has not merely failed at programming or polity. It has contradicted its confession.

That also explains why Principle 12 keeps becoming central. It uses the language of constraint, not suggestion:

> “Every free congregation, as well as every individual believer, is constrained…” 

So yes: **confession is exactly the thread.** The earlier riff was saying the AFLC already has the third way written into its principles. Your current section is pushing the harder consequence: if P1 is confessed, then every member stands under its weight onboard and ashore.



Either institutional authority, where duty is enforced by office and structure, or individual voluntarism, where participation becomes a matter of personal preference, passion, or emotional commitment.

The AFLC principles are more dangerous than either. They say the congregation is free, but not optional. The individual believer is free, but not autonomous. The congregation has no higher earthly ecclesiastical government over it, yet it is still constrained by Word, Spirit, means of grace, Christian love, and mission.
